\clearpage
\VerdierPhysicalPageAnchor{0243}%
\begingroup\small\setlength{\parskip}{0pt}\setlength{\emergencystretch}{1.5em}
\setlength{\abovedisplayskip}{4pt plus 2pt minus 2pt}%
\setlength{\belowdisplayskip}{4pt plus 2pt minus 2pt}%

\VerdierSemanticAnchor{ch4d-4-5-4-proof-a-bcd}
\noindent $D^b(\mathcal A)$ 的对象。态射 $m$ 诱导函子 $F_1$ 与 $F_2$
关于复形 $X$ 的第二典范谱序列之间的态射。在初始项之间诱导的态射
\[
m_2^{p,q}(X)\colon F_1^p\bigl(H^q(X)\bigr)
\longrightarrow F_2^p\bigl(H^q(X)\bigr)
\]
依假设是同构。此外，两个谱序列都是定常的 (4.5.3)。由此可知（第二章，
4.4.5），态射 $m$ 诱导两个谱序列的收敛目标之间的同构，因而 $m(X)$ 是
同构。断言 b)、c)、d) 的证明类似：只须注意，在所作假设下，第二典范谱序列
都是定常的 (4.5.3)。

\VerdierSemanticAnchor{ch4d-4-5-4-proof-e}
证明 e)。根据 c)，只须证明 $\mathcal A$ 的每个对象都是
$\operatorname{Is}(m)$ 的对象。设 $X$ 是 $\mathcal A$ 的对象，$p$ 是整数，
并证明态射
\[
m(X[p])\colon F_1^p(X)\longrightarrow F_2^p(X)
\]
是同构。函子 $F_1$ 与 $F_2$ 左定常，所以存在整数 $r$，使得对
$i=1,2$、$\mathcal A$ 的每个对象 $Y$ 以及每个整数 $q\geq p+r-1$，都有
$F_i^q(Y)=0$。性质 $\alpha$) 蕴涵 $X$ 拟同构于复形

\VerdierSemanticAnchor{ch4d-4-5-4-proof-e-complex-w}
\[
W^\bullet=
\cdots\longrightarrow0\longrightarrow Y\longrightarrow L^{-r+1}
\longrightarrow L^{-r+2}\longrightarrow\cdots\longrightarrow L^{-1}
\longrightarrow L^0\longrightarrow0\longrightarrow\cdots,
\]
其中分量 $L^i$，$-r+1\leq i\leq0$，都是 $M$ 的元素。以 $L^\bullet$
表示复形

\VerdierSemanticAnchor{ch4d-4-5-4-proof-e-complex-l}
\[
\cdots\longrightarrow0\longrightarrow L^{-r+1}
\longrightarrow L^{-r+2}\longrightarrow\cdots\longrightarrow L^{-1}
\longrightarrow L^0\longrightarrow0\longrightarrow\cdots.
\]
有复形的正合列

\VerdierSemanticAnchor{ch4d-4-5-4-proof-e-exact-sequence}
\[
0\longrightarrow L^\bullet\longrightarrow W^\bullet
\longrightarrow Y[r]\longrightarrow0,
\]
从而由 (1.3.2) 得到 $D^*(\mathcal A)$ 中的特异三角：

\VerdierSemanticAnchor{ch4d-4-5-4-proof-e-triangle}
\[
\begin{tikzcd}[column sep=5.2em,row sep=2.8em]
{} & Y[r] \arrow[dl] \\
L^\bullet \arrow[r] & X \arrow[u]
\end{tikzcd}
\]
\endgroup
